% GENERATED FILE. Edit canonical lessons, then run npm run generate:books.
% canonical-source-hash: fnv1a64:aaae77093152ead1
% canonical-lessons: ML-C419-tirannetuppu, ML-C419-vottu, ML-C419-vottar, ML-C419-sthanartthi, ML-C419-balarru

\chapter{Election Day}
\label{ch:ml-election-day}
\hlchaptermodality{419}

\begin{chapteropening}
\textbf{By the end of this chapter:} \emph{I can say an election, a vote, a voter, a candidate and a ballot.}

Complete the last lesson of chapter 419: I can say an election, a vote, a voter, a candidate and a ballot.
\end{chapteropening}

\section[\texorpdfstring{\ml{തിരഞ്ഞെടുപ്പ്} (tiraññeṭuppŭ)}{തിരഞ്ഞെടുപ്പ് (tiraññeṭuppŭ)}]{\ml{തിരഞ്ഞെടുപ്പ്} (tiraññeṭuppŭ) — an election}

\label{lesson:ML-C419-tirannetuppu}



\begin{quote}
Before the new one: say the Malayalam for to prohibit, then the Malayalam for to allow.
\end{quote}

\subsection*{You'll want to know: \ml{തിരഞ്ഞെടുപ്പ്}}
\textbf{\ml{തിരഞ്ഞെടുപ്പ്}} — \emph{tiraññeṭuppŭ} — \textquotedblleft{}an election\textquotedblright{}.

\begin{quote}
\ml{അടുത്ത} \ml{മാസം} \ml{ഞങ്ങളുടെ} \ml{നാട്ടിൽ} \ml{തിരഞ്ഞെടുപ്പ്} \ml{ആണ്}.
\end{quote}

\begin{quote}
\emph{aṭutta māsaṁ ñaṅṅaḷuṭe nāṭṭil tiraññeṭuppŭ āṇŭ.}
\end{quote}

\begin{quote}
Next month there is an election in our area.
\end{quote}

The verb \ml{തിരഞ്ഞെടുക്കുക}, to choose, is the same word.

\begin{cousinweb}
\textbf{\ml{തിരഞ്ഞെടുപ്പ്}} (\emph{tiraññeṭuppŭ}) is a native Malayalam compound: \ml{തിരയുക}, \textquotedblleft{}to seek\textquotedblright{} (DEDR 3242), and \ml{എടുക്കുക}, \textquotedblleft{}to take up, to choose\textquotedblright{} (DEDR 851), with the noun ending -\ml{പ്പ്}. An election is a seeking out and taking up.
\end{cousinweb}

\begin{grammarlens}[title={What you've built}]
One more word for giving reasons and opinions.
\end{grammarlens}

\subsection*{Your turn}
\begin{itemize}
\raggedright
  \item \emph{Say it:} \emph{tiraññeṭuppŭ}
  \item \emph{Say it:} \emph{tiraññeṭuppŭ}, once more
  \item \emph{Say it:} \emph{aṭutta māsaṁ ñaṅṅaḷuṭe nāṭṭil tiraññeṭuppŭ āṇŭ.} — Next month there is an election in our area.
  \item \emph{Recall:} say the Malayalam for to prohibit, then the Malayalam for to allow, then say \emph{tiraññeṭuppŭ} again
\end{itemize}

\begin{tcolorbox}[breakable,colback=teal!4,colframe=teal!35!black,title={Before you move on}]
What does \ml{തിരഞ്ഞെടുപ്പ്} mean? (\textquotedblleft{}An election\textquotedblright{}.) Say it once more.
\end{tcolorbox}
\section[\texorpdfstring{\ml{വോട്ട്} (vōṭṭŭ)}{വോട്ട് (vōṭṭŭ)}]{\ml{വോട്ട്} (vōṭṭŭ) — a vote}

\label{lesson:ML-C419-vottu}



\begin{quote}
Before the new one: say the Malayalam for to allow, then the Malayalam for an election.
\end{quote}

\subsection*{You'll want to know: \ml{വോട്ട്}}
\textbf{\ml{വോട്ട്}} — \emph{vōṭṭŭ} — \textquotedblleft{}a vote\textquotedblright{}.

\begin{quote}
\ml{എന്റെ} \ml{വോട്ട്} \ml{വളരെ} \ml{പ്രധാനമാണ്}.
\end{quote}

\begin{quote}
\emph{enṟe vōṭṭŭ vaḷare pradhānamāṇŭ.}
\end{quote}

\begin{quote}
My vote is very important.
\end{quote}

Plural: \ml{വോട്ടുകൾ}.

\begin{cousinweb}
\textbf{\ml{വോട്ട്}} (\emph{vōṭṭŭ}) was borrowed from English \emph{vote}.
\end{cousinweb}

\begin{grammarlens}[title={What you've built}]
One more word for giving reasons and opinions.
\end{grammarlens}

\subsection*{Your turn}
\begin{itemize}
\raggedright
  \item \emph{Say it:} \emph{vōṭṭŭ}
  \item \emph{Say it:} \emph{vōṭṭŭ}, once more
  \item \emph{Say it:} \emph{enṟe vōṭṭŭ vaḷare pradhānamāṇŭ.} — My vote is very important.
  \item \emph{Recall:} say the Malayalam for to allow, then the Malayalam for an election, then say \emph{vōṭṭŭ} again
\end{itemize}

\begin{tcolorbox}[breakable,colback=teal!4,colframe=teal!35!black,title={Before you move on}]
What does \ml{വോട്ട്} mean? (\textquotedblleft{}A vote\textquotedblright{}.) Say it once more.
\end{tcolorbox}
\section[\texorpdfstring{\ml{വോട്ടർ} (vōṭṭar)}{വോട്ടർ (vōṭṭar)}]{\ml{വോട്ടർ} (vōṭṭar) — a voter}

\label{lesson:ML-C419-vottar}



\begin{quote}
Before the new one: say the Malayalam for an election, then the Malayalam for a vote.
\end{quote}

\subsection*{You'll want to know: \ml{വോട്ടർ}}
\textbf{\ml{വോട്ടർ}} — \emph{vōṭṭar} — \textquotedblleft{}a voter\textquotedblright{}.

\begin{quote}
\ml{ആ} \ml{ഗ്രാമത്തിൽ} \ml{ആയിരം} \ml{വോട്ടർമാർ} \ml{ഉണ്ട്}.
\end{quote}

\begin{quote}
\emph{ā grāmattil āyiraṁ vōṭṭarmār uṇṭŭ.}
\end{quote}

\begin{quote}
There are a thousand voters in that village.
\end{quote}

Plural: \ml{വോട്ടർമാർ}, \emph{vōṭṭarmār}.

Plural with -\ml{മാർ}, the plural for people.

\begin{cousinweb}
\textbf{\ml{വോട്ടർ}} (\emph{vōṭṭar}) was borrowed from English \emph{voter}. The plural takes the Malayalam -\ml{മാർ}.
\end{cousinweb}

\begin{grammarlens}[title={What you've built}]
One more word for giving reasons and opinions.
\end{grammarlens}

\subsection*{Your turn}
\begin{itemize}
\raggedright
  \item \emph{Say it:} \emph{vōṭṭar}
  \item \emph{Say it:} \emph{vōṭṭar}, once more
  \item \emph{Say it:} \emph{ā grāmattil āyiraṁ vōṭṭarmār uṇṭŭ.} — There are a thousand voters in that village.
  \item \emph{Recall:} say the Malayalam for an election, then the Malayalam for a vote, then say \emph{vōṭṭar} again
\end{itemize}

\begin{tcolorbox}[breakable,colback=teal!4,colframe=teal!35!black,title={Before you move on}]
What does \ml{വോട്ടർ} mean? (\textquotedblleft{}A voter\textquotedblright{}.) Say it once more.
\end{tcolorbox}
\section[\texorpdfstring{\ml{സ്ഥാനാർത്ഥി} (sthānārtthi)}{സ്ഥാനാർത്ഥി (sthānārtthi)}]{\ml{സ്ഥാനാർത്ഥി} (sthānārtthi) — a candidate}

\label{lesson:ML-C419-sthanartthi}



\begin{quote}
Before the new one: say the Malayalam for a vote, then the Malayalam for a voter.
\end{quote}

\subsection*{You'll want to know: \ml{സ്ഥാനാർത്ഥി}}
\textbf{\ml{സ്ഥാനാർത്ഥി}} — \emph{sthānārtthi} — \textquotedblleft{}a candidate\textquotedblright{}.

\begin{quote}
\ml{ഇത്തവണ} \ml{മൂന്ന്} \ml{സ്ഥാനാർത്ഥികൾ} \ml{മത്സരിക്കുന്നു}.
\end{quote}

\begin{quote}
\emph{ittavaṇa mūnnŭ sthānārtthikaḷ matsarikkunnu.}
\end{quote}

\begin{quote}
This time three candidates are standing.
\end{quote}

Plural: \ml{സ്ഥാനാർത്ഥികൾ}, \emph{sthānārtthikaḷ}.

Plural: \ml{സ്ഥാനാർത്ഥികൾ}.

\begin{cousinweb}
\textbf{\ml{സ്ഥാനാർത്ഥി}} (\emph{sthānārtthi}) is built from Sanskrit \emph{sthāna}, \textquotedblleft{}a place, a position\textquotedblright{} (CDIAL 13753), and \emph{arthin}, \textquotedblleft{}one who wants, a petitioner\textquotedblright{} (CDIAL 640): one who seeks a post.
\end{cousinweb}

\begin{grammarlens}[title={What you've built}]
One more word for giving reasons and opinions.
\end{grammarlens}

\subsection*{Your turn}
\begin{itemize}
\raggedright
  \item \emph{Say it:} \emph{sthānārtthi}
  \item \emph{Say it:} \emph{sthānārtthi}, once more
  \item \emph{Say it:} \emph{ittavaṇa mūnnŭ sthānārtthikaḷ matsarikkunnu.} — This time three candidates are standing.
  \item \emph{Recall:} say the Malayalam for a vote, then the Malayalam for a voter, then say \emph{sthānārtthi} again
\end{itemize}

\begin{tcolorbox}[breakable,colback=teal!4,colframe=teal!35!black,title={Before you move on}]
What does \ml{സ്ഥാനാർത്ഥി} mean? (\textquotedblleft{}A candidate\textquotedblright{}.) Say it once more.
\end{tcolorbox}
\section[\texorpdfstring{\ml{ബാലറ്റ്} (bālaṟṟŭ)}{ബാലറ്റ് (bālaṟṟŭ)}]{\ml{ബാലറ്റ്} (bālaṟṟŭ) — a ballot}

\label{lesson:ML-C419-balarru}



\begin{quote}
Before the new one: say the Malayalam for a voter, then the Malayalam for a candidate.
\end{quote}

\subsection*{You'll want to know: \ml{ബാലറ്റ്}}
\textbf{\ml{ബാലറ്റ്}} — \emph{bālaṟṟŭ} — \textquotedblleft{}a ballot\textquotedblright{}.

\begin{quote}
\ml{ബാലറ്റ്} \ml{പെട്ടികൾ} \ml{വൈകുന്നേരം} \ml{ഓഫീസിൽ} \ml{എത്തി}.
\end{quote}

\begin{quote}
\emph{bālaṟṟŭ peṭṭikaḷ vaikunnēraṁ ōphīsil etti.}
\end{quote}

\begin{quote}
The ballot boxes reached the office in the evening.
\end{quote}

A ballot box is \ml{ബാലറ്റ്} \ml{പെട്ടി}.

\begin{cousinweb}
\textbf{\ml{ബാലറ്റ്}} (\emph{bālaṟṟŭ}) was borrowed from English \emph{ballot}.
\end{cousinweb}

\begin{grammarlens}[title={What you've built}]
One more word for giving reasons and opinions.
\end{grammarlens}

\subsection*{Your turn}
\begin{itemize}
\raggedright
  \item \emph{Say it:} \emph{bālaṟṟŭ}
  \item \emph{Say it:} \emph{bālaṟṟŭ}, once more
  \item \emph{Say it:} \emph{bālaṟṟŭ peṭṭikaḷ vaikunnēraṁ ōphīsil etti.} — The ballot boxes reached the office in the evening.
  \item \emph{Recall:} say the Malayalam for a voter, then the Malayalam for a candidate, then say \emph{bālaṟṟŭ} again
\end{itemize}

\begin{tcolorbox}[breakable,colback=teal!4,colframe=teal!35!black,title={Before you move on}]
What does \ml{ബാലറ്റ്} mean? (\textquotedblleft{}A ballot\textquotedblright{}.) Say it once more.
\end{tcolorbox}
